\section{Discussion and limitations}
\subsection{What the benchmark result establishes}
The archived 28 October ranking shows that the submitted PHASE-TTC system performs well on GarlTTC Test. Its fusion head operates on phase, timing, and support features whose meanings remain comparable across the two modalities. This avoids requiring independently trained hidden coordinates to align. The ranking alone cannot determine how much of the gain comes from that interface, the longer history, or the upstream models; those questions require matched comparisons. This score concerns the original TRAIN40 graph, not the later quality adapter or canonical-cadence candidate. \evidence{M22,M23,N01--N03}

The out-of-fold training procedure addresses a narrower problem: later stages should not train on features produced by a model fitted to the same sequence group. At deployment, however, they receive features from producers trained on all 40 groups rather than the 24 used in each auxiliary fit. That shift remains, as does the exposure involved in earlier architecture selection. The JEPA experiments provide a separate result: they did not establish a benefit from predictive pretraining, even when embeddings retained variance and effective rank. Subsequent supervised or teacher-assisted gains do not resolve that question.

\subsection{Where fusion helps, and where it fails}
The labeled transfer experiment gives a mixed answer about adding events. Fusion has a lower empirical p95 absolute error than RGB alone, but a higher mean and more errors above ten seconds. All three paired intervals for mean, median, and p95 differences include zero. \evidence{L01,L02,L30,R01}

The range experiments reinforce this point. Improving imminent-contact predictions sometimes worsened estimates for receding objects, because phase error, error in seconds, and warning behavior emphasize different parts of the problem. The labeled transfer includes only positive TTC and cannot test the full signed domain or a complete warning policy.

The original 434-query competition experiment answers a different question. It establishes prediction coverage and disagreement between models under the label status recorded on 28 October. Its Garl visual-only checkpoint also differs from PHASE-TTC's RGB context model. Matching scenario names is insufficient to combine these results with the labeled transfer; the recordings and query identities must match as well. The later corrected target audit enables a separate original-versus-adapted PHASE-TTC error comparison on all 434 queries. It does not change the meaning of the historical disagreement plots or reproduce the published three-sequence Garl protocol. \evidence{T01,T02,N01}

\subsection{Adaptation, cadence, and intermediate objectives}
The quality adapter produces the strongest recent local improvement, reducing query-weighted MAE from 0.824976 to 0.335899\,s under the corrected evaluator. The road-only reduction, from 0.305401 to 0.185200\,s over 365 queries, shows that the aggregate gain is not confined to Slider-750. However, earlier startup and model-selection experiments had inspected the competition recordings. Excluding them from the final 2,608-query fit cannot restore an untouched test set. The adapted predictions also have no verified new official score in this evidence record. \evidence{N01}

Canonical PAIR timing tests a narrower hypothesis. It preserves availability and improves Slider-500, while the descriptive road MAE interval includes zero. The frozen router fails its fixed road screen, and a subsequent 48-fit comparison also fails the combined development criteria. The larger-population aggregate does not revoke that failure. This distinction matters: a correction consistent with units or ideal geometry still has to improve the complete learned estimator under the prespecified controls. \evidence{N02,N03}

The physical and representation experiments reinforce that distinction. Lower future-state loss does not establish a more useful log-height-change representation; algebraic observability does not guarantee a reliable normal-flow estimate; and ideal cuboid inversion does not guarantee better depth from predicted geometry. These negative results delimit the tested mechanisms without proving that predictive learning or physical auxiliaries can never help. Their methodological details remain part of the technical supplement. \evidence{N04--N06}

\subsection{Generalization and deployment}
Most development comparisons reuse sequences or upstream models. Repeating the head with different seeds measures variation conditional on those models, and nearby windows remain correlated observations of the same scene. Independent retraining and new sequences would test how much of the observed performance survives beyond these choices. Deployment also requires obtaining object regions causally; the current evaluations supply them, and rotation-only projection between sensors approximates alignment without depth.

The compact head has 1,955 parameters, but the full system includes all upstream components. A complete parameter and peak-memory inventory remains to be reported. The paced-run records specify the GPU, software, execution settings, and initialization excluded from timing, while omitting the CPU model and GPU driver version. Learned phase widths and propagated variances likewise need empirical calibration before they can support probability statements.

The runtime study identifies an operational bottleneck: all predictions remain numerically consistent, yet their p95 response exceeds 20\,ms. Input preparation and scheduling must improve along with kernel speed. A longer run with physical-camera input and multiple objects would test whether those improvements hold under sustained load. \evidence{L40--L42}

Native cadence replay introduces a further reproducibility boundary. The recaptured original is exactly repeatable on its current platform, yet historical parity fails for intermediate states in 5,805 queries and for TTC in 504. Small original-versus-canonical output changes on those queries measure local sensitivity, not conservation of Test12 accuracy or restoration of historical numerical identity. Private Test12 targets are required to score the new candidates. \evidence{N03}

\section{Conclusion}
PHASE-TTC combines geometric measurements and recent history from events and RGB images, then corrects the RGB estimate through a compact phase-based head. It ranks first in the archived public GarlTTC Test leaderboard of 28 October 2026, reducing Overall MiD by \OfficialOverallRelativeGain{} relative to Garl-TTC. The original labeled transfer still shows a lower p95 alongside worse extreme errors, and the paced service still exceeds its deadline.

The expanded evidence adds a useful but historically exposed quality adaptation on corrected local EvTTC targets. It also shows why temporal normalization, predictive objectives, and physical consistency require complete-system controls: several plausible interventions failed those controls and were not promoted. Independent evaluation of the delivered adapter, resolution of Slider-750's scorer discrepancy, better rare-error behavior, and sustained response deadlines remain central requirements for stronger generalization and deployment claims.
\FloatBarrier
